%% kapa: exam=2 round=24 year=2013 period=1 subject=appraisal_practice subject_ko=감정평가실무 seq=3 number=3 points=20
%% source: Q-Net 감정평가사 기출문제 (cst00309) artlSeq=5077404 file=extracted/1008/5077404/2013년 제24회 감정평가사 2차시험 기출문제/1교시.hwp
%% conversion: kordoc 4.24.51 -> tools/gen2.py (auto); review=proofread-sonnet; flags=none
\begin{longquestion}{【문제 3】}
종전에 시행된 재개발사업으로 인하여 현재 공원으로 이용 중인 토지가 있다. 소유자는 그 사실을 뒤늦게 발견하고 지방자치단체에 보상을 요청하였다. 지방자치단체는 2013. 4. 1을 계약체결일로 하여 보상을 실시하려고 감정평가를 요청하였다. 자료를 참고하여 다음 질문에 답하시오. (20점) \par
\medskip\noindent 물음 1) 대상토지 평가시 적용할 비교표준지의 선정 사유를 설명하시오. (5점) \par
\medskip\noindent 물음 2) 대상토지의 평가시 고려하여야 하는 지목, 실제용도, 지형, 지세, 면적 등을 정리하시오. (5점) \par
\medskip\noindent 물음 3) 실무에서는 개별요인비교치를 가로조건, 접근조건, 환경조건, 획지조건, 행정적조건, 기타조건 등으로 구분하여 산정하고 있다. 자료가 주어지지 않은 조건은 대등한 것으로 보고, 표준지와 대상토지간의 개별요인 비교치를 실무와 같이 산정하고 산정사유를 설명하시오. (5점) \par
\medskip\noindent 물음 4) 제시된 자료에 의해 보상감정평가액을 산정하시오. (5점) \par
\textbf{<자료 1> 종전사업의 개요} \par
1. 사업의 종류: 관악지구 재개발사업 \par
2. 사업인정고시일: 2001. 1. 1 \par
\textbf{<자료 2> 토지조서} \par
{\small\centering
\begin{tabular}{|c|c|c|c|c|c|}
\hline
\multirow{2}{*}{\makecell{기호}} & \multirow{2}{*}{\makecell{소재지}} & \multicolumn{2}{|c|}{\makecell{면적(m\textsuperscript{2})}} & \multirow{2}{*}{\makecell{지목}} & \multirow{2}{*}{\makecell{비고}} \\
\cline{3-4}
 &  & \makecell{공부} & \makecell{편입} &  &  \\
\hline
\makecell{1} & \makecell{관악동 00} & \makecell{1,000} & \makecell{1,000} & \makecell{임야} &  \\
\hline
\end{tabular}\par}

\textbf{<자료 3> 종전사업에 편입될 당시의 이용상황} \par
1. 의뢰된 토지는 종전의 공익사업에 편입되기 전에는 지목이 임야이고, 1970년경부터 무허가건물부지로 이용되던 10,000m\textsuperscript{2}의 토지의 일부였으나, 종전의 공익사업으로 9,000m\textsuperscript{2}는 아파트부지와 도로로 이용 중이고 1,000m\textsuperscript{2}는 공원으로 이용 중이며, 의뢰된 토지는 공원으로 이용 중인 부분임 \par
2. 의뢰된 토지는 현재는 중로에 접하는 장방형의 토지이나, 이는 종전의 공익사업의 시행으로 새로운 도로가 개설되었기 때문이며, 분할되기 이전의 토지는 완만한 경사를 이루고 있는 부정형으로 세로에 접하는 토지였음 \par
\qfigure{figures/03-1.png}
3. 종전의 공익사업이 시행되기 이전에는 일반주거지역이었으나, 종전의 공익사업으로 인하여 대상토지만 자연녹지지역으로 용도지역이 변경되었음 \par
4. 대상토지 주변은 당시 미개발지대로서 남측 근거리의 불량주택지대를 포함하여 재개발사업을 시행하였음 \par
\textbf{<자료 4> 인근의 공시지가 표준지 현황} \par
1. 대상 토지 인근지역에 소재하며 표준적인 이용상황의 표준지는 다음과 같음 \par
{\small\centering
\begin{tabular}{|c|c|c|c|c|c|c|c|c|c|}
\hline
\multirow{2}{*}{\makecell{일련\newline{}번호}} & \multirow{2}{*}{\makecell{소재지}} & \multirow{2}{*}{\makecell{면적\newline{}(m\textsuperscript{2})}} & \multirow{2}{*}{\makecell{지목}} & \multirow{2}{*}{\makecell{이용\newline{}상황}} & \multirow{2}{*}{\makecell{용도\newline{}지역}} & \multirow{2}{*}{\makecell{도로\newline{}교통}} & \multirow{2}{*}{\makecell{형상\newline{}지세}} & \multicolumn{2}{|c|}{\makecell{공시지가(원/m\textsuperscript{2})}} \\
\cline{9-10}
 &  &  &  &  &  &  &  & \makecell{2001년} & \makecell{2013년} \\
\hline
\makecell{1} & \makecell{관악동\newline{}201} & \makecell{1,000} & \makecell{대} & \makecell{단독} & \makecell{일반\newline{}주거} & \makecell{소로\newline{}한면} & \makecell{장방형\newline{}평지} & \makecell{600,000} & \makecell{1,000,000} \\
\hline
\makecell{2} & \makecell{관악동\newline{}202} & \makecell{1,000} & \makecell{대} & \makecell{단독} & \makecell{자연\newline{}녹지} & \makecell{소로\newline{}한면} & \makecell{장방형\newline{}평지} & \makecell{200,000} & \makecell{300,000} \\
\hline
\makecell{3} & \makecell{관악동\newline{}산1} & \makecell{10,000} & \makecell{임야} & \makecell{자연림} & \makecell{일반\newline{}주거} & \makecell{세로\newline{}(가)} & \makecell{부정형\newline{}완경사} & \makecell{120,000} & \makecell{200,000} \\
\hline
\makecell{4} & \makecell{관악동\newline{}산2} & \makecell{10,000} & \makecell{임야} & \makecell{자연림} & \makecell{자연\newline{}녹지} & \makecell{세로\newline{}(가)} & \makecell{부정형\newline{}완경사} & \makecell{20,000} & \makecell{30,000} \\
\hline
\end{tabular}\par}

2. 표준지 일련번호 1은 2000년경 구획정리사업(환지방식)으로 개발된 주거지대 내에 소재하며, 당시 감보율이 40\%였던 것으로 조사되었음 \par
\textbf{<자료 5> 시점 수정자료} \par
1. 2001. 1. 1에서 2013. 4. 1까지의 지가변동률은 60\%임 \par
2. 2013. 1. 1에서 2013. 4. 1까지의 지가변동률은 0\%임 \par
3. 2000년 12월에서 2013년 3월까지의 생산자물가지수 변동률은 40\%임 \par
4. 2012년 12월에서 2013년 3월까지의 생산자물가지수 변동률은 1\%임 \par
\textbf{<자료 6> 각종 격차율} \par
1. 9,000∼10,000m\textsuperscript{2}의 대지는 1,000m\textsuperscript{2} 대지의 85\% 수준임 \par
2. 부정형의 토지는 장방형 토지의 95\% 수준임 \par
3. 대지인 경우 평지는 완경사지의 110\% 수준임 \par
4. 중로에 접하는 토지는 세로에 접하는 토지의 120\% 수준임 \par
5. 소로에 접하는 토지는 세로에 접하는 토지의 105\% 수준임 \par
\textbf{<자료 7> 기타} \par
인근의 평가선례, 매매사례 등과 비교시 비교표준지의 공시지가는 인근지가수준과 같은 수준으로 판단됨 \par
\end{longquestion}
